Z-TRANSFORMS AND BLOCK DIAGRAMS

Z-transform (bilateral)

$$X(z)=\sum_{n=-\infty}^{\infty}x[n]z^{-n}$$

For causal sequences, sum from n=0 upward.

Inverse Z-transform: find sequence from X(z), usually using known pairs, partial fractions, or power series.

BASIC Z-TRANSFORM PAIRS

$$\delta[n]\Leftrightarrow 1$$

$$\delta[n-k]\Leftrightarrow z^{-k}$$

$$a^n u[n]\Leftrightarrow\frac{1}{1-az^{-1}}$$

The causal exponential has ROC $|z|>|a|$.

$$a^{n-k}u[n-k]$$

$$\Leftrightarrow \frac{z^{-k}}{1-az^{-1}}$$

$$a^nu[n-k]\Leftrightarrow\frac{a^kz^{-k}}{1-az^{-1}}$$

Z-TRANSFORM PROPERTIES

$$x[n-k]\Leftrightarrow z^{-k}X(z)$$

$$Ax[n]+By[n]\Leftrightarrow AX(z)+BY(z)$$

$$x[n]*h[n]\Leftrightarrow X(z)H(z)$$

Here * means linear convolution.

For a stable LTI system, evaluate on the unit circle:

$$H(e^{j\omega})=H(z)|_{z=e^{j\omega}}$$

A DELAY IN A BLOCK DIAGRAM

$$z^{-1}=\text{one sample delay}$$

$$z^{-k}=\text{k sample delays}$$

An adder adds incoming branches. A multiplier (gain block) multiplies by its tap coefficient.

TRANSFER FUNCTION FROM DIFFERENCE EQUATION

Assume zero initial conditions. Normalize the coefficient of y[n] to 1:

$$y[n]+\sum_{k=1}^{N}a_ky[n-k]$$

$$=\sum_{k=0}^{M}b_kx[n-k]$$

Then

$$H(z)=\frac{Y(z)}{X(z)}$$

$$H(z)=\frac{\sum_{k=0}^{M}b_kz^{-k}}{1+\sum_{k=1}^{N}a_kz^{-k}}$$

MOST IMPORTANT RULE: FLIP DENOMINATOR SIGNS FOR FEEDBACK TAPS!

To draw the diagram, solve for y[n]:

$$y[n]=\sum_{k=0}^{M}b_kx[n-k]$$

$$-\sum_{k=1}^{N}a_ky[n-k]$$

So numerator gains are $b_k$; feedback gains are $-a_k$.

This sign flip happens because the delayed output terms are moved to the right-hand side.

Example:

$$H(z)=\frac{1+2z^{-1}}{1-0.6z^{-1}+0.2z^{-2}}$$

Multiply across:

$$Y(1-0.6z^{-1}+0.2z^{-2})$$

$$=X(1+2z^{-1})$$

Inverse Z-transform:

$$y[n]-0.6y[n-1]+0.2y[n-2]$$

$$=x[n]+2x[n-1]$$

Solve for output:

$$y[n]=x[n]+2x[n-1]$$

$$+0.6y[n-1]-0.2y[n-2]$$

DRAW ITS DIRECT-FORM I DIAGRAM

1. Input path: x[n] through delay blocks $z^{-1}$; sum gains +1 and +2 into the output adder.

2. Output feedback: branch y[n] through one delay for y[n-1], gain +0.6; two delays for y[n-2], gain -0.2.

3. All four weighted terms enter the SAME adder to produce y[n].

The denominator coefficients are [1, -0.6, +0.2]. The feedback branch gains are [+0.6, -0.2]. Do NOT flip the numerator signs.

DIRECT-FORM II (SHARED DELAYS)

$$w[n]=x[n]-\sum_{k=1}^{N}a_kw[n-k]$$

$$y[n]=\sum_{k=0}^{M}b_kw[n-k]$$

Uses a shared delay line with order max(M,N), instead of separate input and output delay chains.

FROM DIAGRAM TO H(z)

1. Each delay is z^-1.

2. Write the adder's difference equation, including the ACTUAL signs on feedback branches.

3. Transform each delay: x[n-k] becomes z^-k X(z).

4. Collect all Y(z) terms on the left, all X(z) terms on the right, and divide Y(z)/X(z).

BLOCK CONNECTIONS

Cascade (series):

$$H(z)=H_1(z)H_2(z)$$

Parallel (sum):

$$H(z)=H_1(z)+H_2(z)$$

Negative feedback with forward G(z) and return F(z):

$$H(z)=\frac{G(z)}{1+G(z)F(z)}$$

Positive feedback:

$$H(z)=\frac{G(z)}{1-G(z)F(z)}$$

Do not confuse the sign at a feedback summing junction with the sign-flip rule when rearranging a denominator into a difference equation.

POLES, ZEROS, STABILITY

Zeros: roots of the numerator of H(z).

Poles: roots of the denominator of H(z).

For a causal rational LTI filter, BIBO stability requires all actual poles inside the unit circle:

$$|p_k|<1$$

For FIR, the output only uses a finite number of input samples; for IIR, feedback usually produces an infinite impulse response.
